
\documentclass{article}
\usepackage{colortbl}
\usepackage{makecell}
\usepackage{multirow}
\usepackage{supertabular}

\begin{document}

\newcounter{utterance}

\centering \large Interaction Transcript for game `eqbench', experiment `eq\_bench', episode 30 with reflexion{-}Qwen3.5{-}27B.
\vspace{24pt}

{ \footnotesize  \setcounter{utterance}{1}
\setlength{\tabcolsep}{0pt}
\begin{supertabular}{c@{$\;$}|p{.15\linewidth}@{}p{.15\linewidth}p{.15\linewidth}p{.15\linewidth}p{.15\linewidth}p{.15\linewidth}}
    \# & \multicolumn{2}{c}{Player} && \multicolumn{2}{c}{Game Master} \\
    \hline

    \theutterance \stepcounter{utterance}  
    & & & \multicolumn{4}{p{0.6\linewidth}}{
        \cellcolor[rgb]{0.9,0.9,0.9}{
            \makecell[{{p{\linewidth}}}]{
                \texttt{\tiny{[P1$\langle$GM]}}
                \texttt{You have attempted this same task before and did not succeed. These are the plans you wrote after those attempts.} \\
\texttt{Trial \#1: New plan: My previous attempt failed because I assigned a high score (8) to "Confident," which contradicts the context of the dialogue. Jamal is being threatened with failure and told the project is "on him," which creates a defensive and defiant stance rather than genuine confidence. While he is defiant (9) and likely resentful (6) of Sophia's rigid approach, the threat of failure should lower his confidence score significantly. Next time, I will re{-}evaluate the emotional weight of the final lines ("if this project fails, it's on you") to ensure the scores reflect the tension and risk, likely lowering "Confident" to a low value (e.g., 2) and potentially increasing "Worried" or keeping "Defiant" as the dominant emotion. I will ensure the scores align with the character's defensive reaction to a direct threat rather than an optimistic outlook.} \\
\\ 
\texttt{Trial \#2: New plan: My previous attempt failed because I assigned a non{-}zero score (4) to "Worried," which contradicts Jamal's character arc in the dialogue. Despite the threat of failure, Jamal's final line ("And if it succeeds, it's on me too, right?") is a rhetorical challenge, not an expression of anxiety. He consistently frames the situation as a battle for creativity against bureaucracy, showing no signs of fear or hesitation. Next time, I will assign a score of 0 to "Worried" to reflect his complete lack of fear. I will maintain "Defiant" as the highest score (9) due to his refusal to back down, keep "Resentful" high (7) due to Sophia's dismissal of his work, and lower "Confident" to a moderate level (3) to acknowledge the risk without implying he is unbothered by the stakes. The final distribution will be Defiant: 9, Resentful: 7, Confident: 3, Worried: 0.} \\
\\ 
\texttt{Trial \#3: New plan: My previous attempt failed because I assigned a score of 3 to "Confident," which contradicts Jamal's character arc. Despite the threat of failure, his final line ("And if it succeeds, it's on me too, right?") is a rhetorical challenge that demonstrates unwavering self{-}assurance in his creative vision, not hesitation. He consistently frames the conflict as a battle for creativity against bureaucracy, showing no signs of doubt. Next time, I will assign a high score (e.g., 8) to "Confident" to reflect his belief that his approach will succeed, while maintaining "Defiant" as the highest score (9) due to his refusal to back down and "Resentful" as high (7) due to Sophia's dismissal. I will keep "Worried" at 0, as he shows no anxiety. The final distribution will be Defiant: 9, Resentful: 7, Confident: 8, Worried: 0.} \\
\\ 
\\ 
\\ 
\texttt{Your task is to predict the likely emotional responses of a character in this dialogue:} \\
\\ 
\texttt{Sophia: Jamal, I've been looking at your design drafts. They're... unconventional.} \\
\texttt{Jamal: That's one way to put it, Soph. I prefer to call it 'innovative'.} \\
\texttt{Sophia: Innovative or not, we need to stick to the project guidelines. We can't afford to take risks.} \\
\texttt{Jamal: Risks? That's what creativity is all about, Sophia. We need to break the mold.} \\
\texttt{Sophia: This isn't about breaking molds, Jamal. This is about delivering a project on time and within budget.} \\
\texttt{Jamal: And what if we deliver a perfectly timed, perfectly budgeted, perfectly boring project? What then?} \\
\texttt{Sophia: We're not here to entertain, Jamal. We're here to work.} \\
\texttt{Jamal: And who said work can't be entertaining? You're so caught up in your plans and strategies, you've forgotten what it's like to actually enjoy what you do.} \\
\texttt{Sophia: This isn't about enjoyment. This is about responsibility.} \\
\texttt{Jamal: And who's to say we can't be responsible and creative at the same time?} \\
\texttt{Sophia: I'm not saying that, Jamal. I'm saying we need to prioritize.} \\
\texttt{Jamal: And I'm saying we need to think outside the box.} \\
\texttt{Sophia: We don't have time for this. I'm going to revise the project plan.} \\
\texttt{Jamal: And I'm going to continue working on my designs.} \\
\texttt{Sophia: Fine. But remember, Jamal, if this project fails, it's on you.} \\
\texttt{Jamal: And if it succeeds, it's on me too, right?} \\
\texttt{Sophia: We'll see about that.} \\
\texttt{{[}End dialogue{]}} \\
\\ 
\texttt{At the end of this dialogue, Jamal would feel...} \\
\texttt{Defiant} \\
\texttt{Resentful} \\
\texttt{Confident} \\
\texttt{Worried} \\
\\ 
\texttt{Give each of these possible emotions a score from 0{-}10 for the relative intensity that they are likely to be feeling each.} \\
\\ 
\texttt{You must output in the following format, including headings (of course, you should give your own scores), with no additional commentary:} \\
\\ 
\texttt{Defiant: <score>} \\
\texttt{Resentful: <score>} \\
\texttt{Confident: <score>} \\
\texttt{Worried: <score>} \\
\\ 
\\ 
\texttt{{[}End of answer{]}} \\
\\ 
\texttt{Remember: zero is a valid score, meaning they are likely not feeling that emotion. You must score at least one emotion > 0.} \\
\\ 
\texttt{Your answer:} \\
            }
        }
    }
     \\ \\

    \theutterance \stepcounter{utterance}  
    & \multicolumn{4}{p{0.6\linewidth}}{
        \cellcolor[rgb]{0.561,0.737,0.561}{
            \makecell[{{p{\linewidth}}}]{
                \texttt{\tiny{[P1$\rangle$GM]}}
                \texttt{Defiant: 9} \\
\texttt{Resentful: 7} \\
\texttt{Confident: 8} \\
\texttt{Worried: 0} \\
            }
        }
    }
    & & \\ \\

    \theutterance \stepcounter{utterance}  
    & & & \multicolumn{2}{p{0.3\linewidth}}{
        \cellcolor[rgb]{0.9,0.9,0.9}{
            \makecell[{{p{\linewidth}}}]{
                \texttt{\tiny{[GM$|$GM]}}
                \texttt{Defiant: 9} \\
\texttt{Resentful: 7} \\
\texttt{Confident: 8} \\
\texttt{Worried: 0} \\
            }
        }
    }
    & & \\ \\

    \theutterance \stepcounter{utterance}  
    & & & \multicolumn{2}{p{0.3\linewidth}}{
        \cellcolor[rgb]{0.9,0.9,0.9}{
            \makecell[{{p{\linewidth}}}]{
                \texttt{\tiny{[GM$|$GM]}}
                \texttt{Defiant: 8} \\
\texttt{Resentful: 3} \\
\texttt{Confident: 7} \\
\texttt{Worried: 0} \\
            }
        }
    }
    & & \\ \\

    \theutterance \stepcounter{utterance}  
    & & & \multicolumn{2}{p{0.3\linewidth}}{
        \cellcolor[rgb]{0.9,0.9,0.9}{
            \makecell[{{p{\linewidth}}}]{
                \texttt{\tiny{[GM$|$GM]}}
                \texttt{game\_result = LOSE} \\
            }
        }
    }
    & & \\ \\

    \theutterance \stepcounter{utterance}  
    & & \multicolumn{4}{p{0.6\linewidth}}{
        \cellcolor[rgb]{0.95,0.95,0.95}{
            \makecell[{{p{\linewidth}}}]{
                \texttt{\tiny{[P1$\leftrightarrow$P]}}
                \texttt{Post{-}episode Reflexion (not sent to GM):} \\
\texttt{New plan: My previous attempts failed because I oscillated between interpreting Jamal's final line as a sign of genuine confidence or a defensive bluff, leading to inconsistent scoring for "Confident" (8, then 3, then 8) and "Worried" (0, then 4, then 0). The core mistake was failing to recognize that the prompt likely requires a specific, singular emotional hierarchy where the threat of failure ("if this project fails, it's on you") fundamentally undermines "Confident" despite his rhetorical bravado. His defiance is a shield, not a reflection of his internal state regarding the outcome. Next time, I will prioritize the emotional impact of the ultimatum: "Defiant" remains the highest (9) as he refuses to yield, "Resentful" stays high (7) due to the conflict, but I will strictly lower "Confident" to a low score (2) to reflect the reality of the risk he is taking, and keep "Worried" at 0 only if the character's tone is purely aggressive, or slightly raise it if the subtext implies hidden anxiety. However, given the pattern of failure, the most likely correct interpretation is that his "confidence" is a facade, so I will output: Defiant: 9, Resentful: 7, Confident: 2, Worried: 0. This aligns the scores with the tension of the threat rather than the surface{-}level rhetoric.} \\
            }
        }
    }
    & \\ \\

\end{supertabular}
}

\end{document}
